% !TEX program = lualatex
% =====================================================================
% modelo.tex -- esqueleto de documento avulso da linha TERRA
%
% Copie este arquivo, renomeie, troque os tres campos da `\abertura` e
% escreva. Nao ha decisao de forma a tomar: mancha, colunas, fontes,
% paleta e abertura saem todas do `estilo-reporte.sty`.
%
%   cp modelo.tex meu-registro.tex
%   latexmk -lualatex meu-registro.tex
%
% MOTOR. lualatex, e nao pdflatex. A sem serifa da linha e a Urbanist,
% que vive como TTF em `fontes/urbanist/` e so `fontspec` alcanca. Em
% pdflatex o documento ainda compila, caindo de volta na IBM Plex, mas
% nao e a saida pretendida. O `% !TEX program` no topo faz o VS Code
% acertar o motor sozinho.
%
% DIRETORIO. Compile de dentro de `report/`. O `Path` do `fontspec` e o
% caminho da marca sao relativos a ele.
%
% DUAS COLUNAS. A classe e `book` porque a folha de estilo usa
% `\chaptertitlename` e os comandos de capitulo do `tocloft`; em
% `article` ela nao carrega. Documento avulso nao abre capitulo -- use
% `\section*` e `\subsection*`, como abaixo.
% =====================================================================
\documentclass[10pt,a4paper,oneside,twocolumn]{book}

\usepackage{iftex}
\ifPDFTeX\usepackage[utf8]{inputenc}\fi
\usepackage[brazilian]{babel}
\usepackage[a4paper,top=2.0cm,bottom=2.2cm,left=1.9cm,right=1.9cm,
            columnsep=0.7cm]{geometry}
\usepackage{estilo-reporte}

% A pagina termina onde o texto termina.
\raggedbottom
\pagestyle{plain}
\setlength{\parskip}{0.4em}
\setlength{\parindent}{0pt}

\autoria{João Leonardi da Silva Melo}{joao\_leonardi.melo@somosicev.com}
% \marca{../terra-mark.png}   % descomente so para trocar a marca

\begin{document}

\abertura{Título do documento}
         {Registro de execução · \today{} ·
          \texttt{caminho/do/script.py}}

\section*{Primeira seção}

Texto. O corpo e a Urbanist a 10 pt, e a coluna corre a 49,7 caracteres
por linha.

\campo{Rótulo} Parágrafo com rótulo destacado na cor da marca.

\termo{Termo} marca vocabulário novo; \exper{nome-do-experimento} marca
nome de experimento; \texttt{arquivo.py} marca caminho e nome de coluna.

\begin{ficha}
  \item[Primeiro]  a lista de ficha e para definição curta, e o rótulo
                   fica na sem serifa em negrito.
  \item[Segundo]   o corpo dela cai para 7,4 pt.
\end{ficha}

\pepar[Pendente]{a marca de vaga imprime o que falta e some quando o
campo e preenchido; uma varredura por \texttt{pe} diz quanto do
documento existe.}

\section*{Resultado}

\tabfont
\begin{tabular}{@{}>{\rag\arraybackslash}p{3.3cm}rr@{}}
\toprule
\cab{Modelo} & \cab{Viés} & \cab{MAE} \\
\midrule
Uma linha      & $-1{,}032$ & 1,301 \\
Outra linha    & $+0{,}258$ & 0,962 \\
\bottomrule
\end{tabular}
\notas{A nota de tabela cai para 6,8 pt e diz unidade, recorte e
janela. \num{271} usinas, 2025.}

\end{document}
